%% ============================================================
%% Affective Computing — Week 2 Study Notes
%% ============================================================
\documentclass[11pt,a4paper]{article}

\usepackage[a4paper, left=1.8cm, right=1.8cm, top=1.3cm, bottom=1.8cm]{geometry}
\usepackage{booktabs}
\usepackage{tabularx}
\usepackage{array}
\usepackage{enumitem}
\usepackage{titlesec}
\usepackage{fancyhdr}
\usepackage{fontenc}
\usepackage{inputenc}
\usepackage{microtype}
\usepackage{parskip}
\usepackage{hyperref}
\usepackage{amsmath}
\usepackage{amssymb}

\hypersetup{hidelinks, colorlinks=false}

%% ── Table helpers ─────────────────────────────────────────
\newcolumntype{L}[1]{>{\raggedright\arraybackslash}p{#1}}

%% ── Section headings ──────────────────────────────────────
\titleformat{\section}{\large\bfseries}{}{0em}{}[\titlerule]
\titleformat{\subsection}{\normalsize\bfseries}{}{0em}{}
\titleformat{\subsubsection}{\normalsize\bfseries}{}{0em}{}

\titlespacing{\section}{0pt}{12pt}{4pt}
\titlespacing{\subsection}{0pt}{8pt}{3pt}
\titlespacing{\subsubsection}{0pt}{6pt}{2pt}

%% ── Page style ────────────────────────────────────────────
\pagestyle{fancy}
\fancyhf{}
\renewcommand{\headrulewidth}{0pt}
\fancyfoot[C]{\small Affective Computing --- Week 2 \textbar\ Page \thepage}

%% ── Misc helpers ──────────────────────────────────────────
\newcommand{\bb}[1]{\textbf{#1}}
\setlength{\parskip}{4pt}

%% ═══════════════════════════════════════════════════════════
\begin{document}
\setlength{\arrayrulewidth}{0.4pt}

\begin{center}
\bfseries\LARGE Affective Computing — Week 2\par
\vspace{0.2cm}
\large Assignment 2 Detailed Solutions
\end{center}
\vspace{0.5cm}

%% ════════════════════════════════════════════════════════════
\section*{ASSIGNMENT 2 --- Full Solutions with Explanations}

\subsubsection*{Question 1}
\textit{Which brain region is primarily involved in regulating emotional experience and motivation?}
\medskip

\begin{tabular}{L{0.3cm} L{14.9cm}}
& \textbf{A.\ Prefrontal cortex \checkmark} \\
& B.\ Occipital lobe \\
& C.\ Cerebellum \\
& D.\ Pituitary gland \\
\end{tabular}

\vspace{0.3cm}
\noindent\textbf{Ans:} A --- Prefrontal cortex\\[0.2cm]
\textbf{Why this answer:} \\
The \textbf{prefrontal cortex (PFC)} plays a central role in regulating emotional experience and motivation. It is involved in emotional appraisal, decision-making, and modulating emotional responses generated by subcortical structures like the amygdala.
\vspace{0.4cm}

\subsubsection*{Question 2}
\textit{The amygdala is most strongly associated with which function?}
\medskip

\begin{tabular}{L{0.3cm} L{14.9cm}}
& A.\ Decision making and planning \\
& B.\ Regulation of voluntary movement \\
& \textbf{C.\ Processing stimulus salience and arousal \checkmark} \\
& D.\ Language comprehension \\
\end{tabular}

\vspace{0.3cm}
\noindent\textbf{Ans:} C --- Processing stimulus salience and arousal\\[0.2cm]
\textbf{Why this answer:} \\
The \textbf{amygdala} is primarily responsible for detecting emotionally significant stimuli and processing their \textbf{salience} (relevance) and \textbf{arousal} value. It plays a key role in fear conditioning, threat detection, and tagging stimuli as emotionally important.
\vspace{0.4cm}

\subsubsection*{Question 3}
\textit{According to embodied cognition theories, voluntary facial muscle contraction can influence emotional experience. This idea supports which concept?}
\medskip

\begin{tabular}{L{0.3cm} L{14.9cm}}
& A.\ Emotions exist independently of bodily states \\
& B.\ Emotion cannot occur without a conscious appraisal \\
& C.\ Emotion is generated only through cognitive evaluation \\
& \textbf{D.\ Bodily feedback contributes to how emotions are felt \checkmark} \\
\end{tabular}

\vspace{0.3cm}
\noindent\textbf{Ans:} D --- Bodily feedback contributes to how emotions are felt\\[0.2cm]
\textbf{Why this answer:} \\
Embodied cognition theories propose that \textbf{bodily feedback} (including facial muscle contractions) directly contributes to emotional experience. This aligns with the \textbf{facial feedback hypothesis}, which shows that voluntarily smiling or frowning can modulate the intensity of felt emotions.
\vspace{0.4cm}
\newpage
\subsubsection*{Question 4}
\textit{Higher activation in the left frontal hemisphere is generally linked to which motivational tendency?}
\medskip

\begin{tabular}{L{0.3cm} L{14.9cm}}
& A.\ Withdrawal motivation \\
& \textbf{B.\ Approach motivation \checkmark} \\
& C.\ Reduced emotional awareness \\
& D.\ Neutral affect \\
\end{tabular}

\medskip

\textbf{Why this answer:} \\
Research on \textbf{frontal asymmetry} shows that greater \textbf{left frontal activation} is associated with \textbf{approach motivation} --- the tendency to move toward goals, rewards, and positive stimuli. Conversely, greater right frontal activation is linked to withdrawal motivation.
\vspace{0.4cm}

\subsubsection*{Question 5}
\textit{Which emotion is commonly associated with activation of the left frontal PFC, despite being negatively valenced?}
\medskip

\begin{tabular}{L{0.3cm} L{14.9cm}}
& A.\ Fear \\
& B.\ Sadness \\
& \textbf{C.\ Anger \checkmark} \\
& D.\ Disgust \\
\end{tabular}

\vspace{0.3cm}
\noindent\textbf{Ans:} B --- Approach motivation\\[0.2cm]
\textbf{Why this answer:} \\
\textbf{Anger} is a unique case --- despite being a negatively valenced emotion, it activates the \textbf{left frontal PFC} because it is fundamentally an \textbf{approach-oriented} emotion. Anger motivates action toward confronting an obstacle or threat, aligning it with approach motivation rather than withdrawal.
\vspace{0.4cm}

\subsubsection*{Question 6}
\textit{In the dimensional VAD/PAD model, dominance refers to:}
\medskip

\begin{tabular}{L{0.3cm} L{14.9cm}}
& \textbf{A.\ The sense of control or submission associated with an emotion \checkmark} \\
& B.\ The pleasantness of an emotional state \\
& C.\ The intensity or energy of an emotion \\
& D.\ The clarity with which an emotion is recognized \\
\end{tabular}

\vspace{0.3cm}
\noindent\textbf{Ans:} A --- The sense of control or submission associated with an emotion\\[0.2cm]
\textbf{Why this answer:} \\
In the \textbf{VAD/PAD (Valence-Arousal-Dominance)} model, \textbf{dominance} captures the degree of control or power an individual feels in an emotional state. High dominance (e.g., anger) reflects feeling in control, while low dominance (e.g., fear) reflects feeling submissive or helpless.
\vspace{0.4cm}

\subsubsection*{Question 7}
\textit{One limitation of traditional categorical emotion models is that:}
\medskip

\begin{tabular}{L{0.3cm} L{14.9cm}}
& A.\ They cannot be implemented computationally \\
& \textbf{B.\ They struggle to represent relationships between emotions \checkmark} \\
& C.\ They include too many dimensions to be practical \\
& D.\ They ignore facial expressions entirely \\
\end{tabular}

\vspace{0.3cm}
\noindent\textbf{Ans:} B --- They struggle to represent relationships between emotions\\[0.2cm]
\textbf{Why this answer:} \\
Traditional \textbf{categorical models} (e.g., Ekman's basic emotions) classify emotions into discrete categories but struggle to capture \textbf{relationships, gradations, and overlaps} between emotions. For example, they don't easily represent how irritation transitions into anger, or how emotions can blend together.
\vspace{0.4cm}
\newpage
\subsubsection*{Question 8}
\textit{Which emotion is typically characterized by lowered eyebrows, tense eyelids, and pressed lips?}
\medskip

\begin{tabular}{L{0.3cm} L{14.9cm}}
& A.\ Happiness \\
& B.\ Fear \\
& C.\ Disgust \\
& \textbf{D.\ Anger \checkmark} \\
\end{tabular}

\vspace{0.3cm}
\noindent\textbf{Ans:} D --- Anger\\[0.2cm]
\textbf{Why this answer:} \\
The facial expression of \textbf{anger} is characterized by \textbf{lowered and drawn-together eyebrows}, \textbf{tense/narrowed eyelids}, and \textbf{tightly pressed or thinned lips}. These Action Units (AUs) in the Facial Action Coding System (FACS) distinguish anger from other negative emotions like fear or disgust.
\vspace{0.4cm}

\subsubsection*{Question 9}
\textit{Increased skin conductance is most reliably associated with which type of disgust-related stimulus?}
\medskip

\begin{tabular}{L{0.3cm} L{14.9cm}}
& A.\ Pictures of dirty toilets \\
& B.\ Pleasant odors \\
& \textbf{C.\ Mutilation or body-boundary violation images \checkmark} \\
& D.\ Neutral facial expressions \\
\end{tabular}

\vspace{0.3cm}
\noindent\textbf{Ans:} C --- Mutilation or body-boundary violation images\\[0.2cm]
\textbf{Why this answer:} \\
Research shows that \textbf{mutilation and body-boundary violation images} produce the most reliable increase in \textbf{skin conductance (EDA)} among disgust-related stimuli. These stimuli elicit strong autonomic arousal because they involve threats to bodily integrity, triggering both disgust and a defensive physiological response.
\vspace{0.4cm}

\subsubsection*{Question 10}
\textit{Crying-related sadness is associated with which physiological pattern?}
\medskip

\begin{tabular}{L{0.3cm} L{14.9cm}}
& A.\ Reduced respiration and lower HRV \\
& \textbf{B.\ Increased heart rate and increased skin conductance \checkmark} \\
& C.\ Decreased skin conductance and reduced heart rate \\
& D.\ No change in respiration but reduced HRV \\
\end{tabular}

\vspace{0.3cm}
\noindent\textbf{Ans:} B --- Increased heart rate and increased skin conductance\\[0.2cm]
\textbf{Why this answer:} \\
\textbf{Crying-related sadness} is physiologically distinct from non-crying sadness. It is associated with \textbf{increased heart rate} and \textbf{increased skin conductance}, reflecting heightened sympathetic nervous system activation. This makes crying an active physiological state, unlike passive sadness which shows reduced arousal.
\vspace{0.4cm}

\medskip
\subsection*{Assignment Quick Answer Summary}

\begin{center}
\renewcommand{\arraystretch}{1.4}
\begin{tabular}{|L{0.8cm}|L{6.2cm}|L{8.2cm}|}
\hline
\bb{Q\#} & \bb{Topic} & \bb{Correct Answer} \\
\hline
1 & Brain Regions \& Emotion & A --- Prefrontal cortex \\
2 & Amygdala Function & C --- Processing stimulus salience and arousal \\
3 & Embodied Cognition & D --- Bodily feedback contributes to emotions \\
4 & Frontal Asymmetry & B --- Approach motivation \\
5 & Anger \& Left Frontal PFC & C --- Anger \\
6 & VAD/PAD Model --- Dominance & A --- Sense of control or submission \\
7 & Categorical Model Limitations & B --- Struggle to represent relationships \\
8 & Facial Expression of Anger & D --- Anger \\
9 & Skin Conductance \& Disgust & C --- Mutilation/body-boundary violation images \\
10 & Crying-Related Sadness Physiology & B --- Increased heart rate and skin conductance \\
\hline
\end{tabular}
\end{center}

\medskip
{\small\textit{Reference: Affective Computing --- Week 2 Lecture Materials}}

\end{document}
